\chapter{Sectioning}

KOMA-Script extends the standard LaTeX sectioning commands with additional
features such as flexible numbering, paragraph hooks, and style customisation.
OmniLaTeX further tailours the sectioning defaults for a clean, modern look.

\section{KOMA Section Hierarchy}

The full hierarchy, from coarsest to finest:

\begin{verbatim}
\part          % I, II, III, ...
\chapter       % 1, 2, 3, ...      (book/report only)
\section       % 1.1, 1.2, ...
\subsection    % 1.1.1, 1.1.2, ...
\subsubsection % 1.1.1.1, ...
\paragraph     % 1.1.1.1.1, ...
\subparagraph  % 1.1.1.1.1.1, ...
\end{verbatim}

Each level is numbered by default.  The depth at which numbering and TOC
entries appear is controlled by two counters:

\begin{verbatim}
\setcounter{secnumdepth}{3}  % number up to \subsubsection
\setcounter{tocdepth}{2}     % show up to \subsection in TOC
\end{verbatim}

OmniLaTeX defaults to \texttt{secnumdepth=3} and \texttt{tocdepth=2}.

\section{Chapter Prefix}

OmniLaTeX scales the chapter number to 4.5$\times$ its normal size using
KOMA's font interface:

\begin{verbatim}
\setkomafont{chapter}{\fontsize{45}{54}\selectfont\bfseries}
\end{verbatim}

This creates a prominent visual anchor for each chapter opening.  The
chapter title itself is set in the default sans-serif font at normal size,
creating a strong typographic contrast between the large number and the
readable heading text.

To customise the chapter number size:

\begin{verbatim}
\setkomafont{chapter}{\fontsize{60}{72}\selectfont\bfseries}
\end{verbatim}

\section{Table of Contents}

Generate a TOC with \verb|\tableofcontents|.  OmniLaTeX styles the TOC
with a clean sans-serif layout and dotted leaders:

\begin{verbatim}
\frontmatter
\tableofcontents
\mainmatter
\end{verbatim}

Control which sectioning levels appear:

\begin{verbatim}
\setcounter{tocdepth}{1}  % chapters only
\setcounter{tocdepth}{3}  % down to \subsubsection
\end{verbatim}

Individual entries can be hidden with \verb|\addtocontents{toc}{\vspace{1em}}|
or by using the \texttt{tocdepth} counter locally within a chapter.

\section{Lists of Floats}

KOMA-Script provides built-in lists for the standard float types:

\begin{verbatim}
\listoffigures   % List of Figures
\listoftables    % List of Tables
\end{verbatim}

OmniLaTeX also registers \texttt{example} as a float type, enabling:

\begin{verbatim}
\listofexamples  % List of Examples
\end{verbatim}

For custom float types, declare a new list with \verb|\DeclareNewTOC|:

\begin{verbatim}
\DeclareNewTOC[
  type=algorithm,
  name=List of Algorithms,
  listname={Algorithm}
]{loa}
\end{verbatim}

This creates a \verb|\listofalgorithms| command and a \texttt{loa} file
for the auxiliary data.

\section{PDF Bookmarks}

PDF bookmarks (the navigable outline in PDF viewers) are generated
automatically from sectioning commands.  Add manual bookmarks with:

\begin{verbatim}
\pdfbookmark[0]{Title Page}{titlepage}
\pdfbookmark[1]{Appendix A}{appa}
\end{verbatim}

The optional argument sets the bookmark level (0 = top-level, matching
\verb|\part|; 1 = chapter level).

OmniLaTeX enables bold chapter bookmarks by default via
\verb|\boldPDFchapterstrue|.  To disable:

\begin{verbatim}
\boldPDFchaptersfalse
\end{verbatim}

\section{Unnumbered Chapters}

For chapters that should not appear in the numbering sequence:

\begin{verbatim}
\addchap{Acknowledgements}   % KOMA: unnumbered, appears in TOC
\chapter*{Colophon}          % LaTeX: unnumbered, absent from TOC
\end{verbatim}

\begin{description}
  \item[\texttt{\textbackslash addchap}] --- KOMA-Script command.  The chapter
        is unnumbered but still listed in the TOC and generates a PDF bookmark.
  \item[\texttt{\textbackslash chapter*}] --- Standard LaTeX starred command.
        The chapter is unnumbered and does \emph{not} appear in the TOC or PDF
        bookmarks.
\end{description}

OmniLaTeX recommends \verb|\addchap| for front-matter chapters (abstract,
acknowledgements, nomenclature) that should be listed but not numbered.
